% !TeX root = main.tex
\chapter{Conclusions and Future Work}

\section{Conclusions}

\quad This thesis studied Java-Python program translation on the AVATAR
benchmark, with an emphasis on whether modern pretrained code models can
produce translations that are not only textually similar to the reference, but
also functionally correct. The work compared four model categories: a
zero-shot instruction model, an encoder-only model, a decoder-only model, and a
Seq2Seq encoder-decoder model. This organization made it possible to compare
different ways of using pretraining: prompting a large model directly, reusing
only a source encoder, adapting a causal decoder with parameter-efficient
fine-tuning, and fine-tuning a pretrained encoder-decoder model.

The results show that Computational Accuracy is the most important metric for
this task. BLEU and CodeBLEU are useful because they measure token-level and
code-aware similarity, but they do not guarantee that the translated program
preserves the behavior of the source program. This was visible both in the
original AVATAR paper and in the experiments from this thesis. Several models
produced translations that looked close to the reference, but only a smaller
part of those translations passed the available execution tests.

The strongest implemented system was the zero-shot Qwen2.5-Coder 7B model. It
reached the highest Computational Accuracy in both translation directions,
with 59.68 for Java-to-Python and 61.80 for Python-to-Java, even though it was
not fine-tuned on AVATAR. This result shows the strength of recent
instruction-tuned code models: a sufficiently capable model can solve many
translation examples from prompting alone. The Seq2Seq CodeT5+ 220M model was
the strongest fine-tuned system. It achieved the highest BLEU scores among the
implemented models and strong Computational Accuracy values of 47.82 and
41.29. This confirms that pretrained encoder-decoder models are a strong fit
for supervised code translation.

The other two model categories were useful because they clarified the limits
of partial adaptation. The decoder-only Qwen2.5-Coder 0.5B model improved over
the weaker reference results from the original AVATAR paper, but it remained
below the larger zero-shot model and the Seq2Seq CodeT5+ model. The
encoder-only UniXcoder model produced reasonable lexical similarity, but its
functional correctness remained low. This shows that source-side pretraining
alone is not enough for reliable program translation; the target-side generator
also needs enough capacity and task adaptation.

All experiments were implemented and run in a reproducible project structure.
Kaggle was used as the main execution environment because it provided access to
GPU resources and made it possible to train and evaluate models without a local
high-end machine. The complete code for the project is available at
\url{https://github.com/nicosuhan/avatar}.

\section{Future Work}

\quad The main limitation of this work was computational budget. Because the
experiments were run on limited Kaggle resources, there was not enough time to
train every model for more epochs or to run a wider hyperparameter search.
Future work should therefore repeat the supervised experiments with longer
training schedules, especially for the decoder-only and Seq2Seq models. More
epochs could help the models adapt better to AVATAR's token-spaced format and
could reduce the gap between lexical similarity and functional correctness.

Another direction is to test larger checkpoints under the same evaluation
protocol. The zero-shot 7B model performed very well, while the fine-tuned
0.5B decoder-only model was much weaker. This suggests that model scale has a
large effect on functional translation quality. Future experiments could
fine-tune larger decoder-only models with LoRA or QLoRA
\cite{hu2022lora,dettmers2023qlora} and compare them directly with the
zero-shot setup.

Finally, future work could improve the evaluation and error analysis. The
current results show how many programs pass the available tests, but a deeper
analysis could separate syntax errors, API mismatches, input-output mistakes,
and algorithmic errors. This would make it easier to understand why a model
fails and which architecture is best suited for each type of translation
problem. Together with longer training and larger models, this would provide a
clearer path toward more reliable Java-Python program translation.
